\documentclass[12pt]{article}
\usepackage[utf8]{inputenc}
\usepackage[T1]{fontenc}
\usepackage[margin=1in]{geometry}
\usepackage{setspace}
\usepackage{xcolor}
\usepackage[colorlinks=true, urlcolor=blue!60!black, linkcolor=blue!60!black]{hyperref}
\setlength{\parindent}{2em}
\setlength{\parskip}{1.25em}
\renewcommand{\baselinestretch}{1.0}
\begin{document}
\begin{center}
{\LARGE \textit{Massachusetts v. EPA}, 549 U.S. 497 (2007)\par}
\vspace{0.8em}
{\large Case Brief\par}
\end{center}
\vspace{1em}

\section{Memory Jogger}

\textbf{Case \#164 --- Course Unit: Standing.} In a landmark 5--4 decision, the Supreme Court held that Massachusetts had Article III standing---bolstered by ``special solicitude'' for a sovereign State asserting a statutory procedural right to protect its quasi-sovereign territory---to challenge the Environmental Protection Agency's (EPA) refusal to regulate greenhouse gas emissions from new motor vehicles, even though climate change is a global, widely shared harm and domestic vehicle standards would only incrementally slow it. On the merits, the Court held that greenhouse gases unambiguously qualify as ``air pollutants'' under the Clean Air Act and that EPA's policy-based refusal to make an endangerment judgment was arbitrary and capricious (\href{https://supreme.justia.com/cases/federal/us/549/497/}{Justia}; \href{https://tile.loc.gov/storage-services/service/ll/usrep/usrep549/usrep549497/usrep549497.pdf}{Official U.S. Reports PDF}).

\section{Detailed Case Facts}

On October 20, 1999, nineteen private environmental and renewable-energy organizations petitioned the EPA to initiate rulemaking to regulate emissions of four heat-trapping greenhouse gases---carbon dioxide ($\mathrm{CO_2}$), methane, nitrous oxide, and hydrofluorocarbons---from new motor vehicles under Section 202(a)(1) of the Clean Air Act (CAA), 42 U.S.C. \S\ 7521(a)(1) (\href{https://www.law.cornell.edu/supct/html/05-1120.ZO.html}{Cornell LII, Majority Opinion}). Section 202(a)(1) directs that the EPA Administrator ``shall by regulation prescribe . . . standards applicable to the emission of any air pollutant from any class or classes of new motor vehicles or new motor vehicle engines, which in his judgment cause, or contribute to, air pollution which may reasonably be anticipated to endanger public health or welfare'' (\href{https://supreme.justia.com/cases/federal/us/549/497/}{Justia}). The Act defines ``air pollutant'' to include ``any air pollution agent or combination of such agents, including any physical, chemical, biological, radioactive . . . substance or matter which is emitted into or otherwise enters the ambient air,'' 42 U.S.C. \S\ 7602(g), and defines ``welfare'' to include ``effects on . . . weather . . . and climate,'' \S\ 7602(h) (\href{https://supreme.justia.com/cases/federal/us/549/497/}{Justia}).

Prior to the petition, two consecutive EPA General Counsels (Jonathan Z. Cannon in April 1998 and Gary S. Guzy in October 1999) had concluded that $\mathrm{CO_2}$ fell within EPA's regulatory authority under the CAA, though the agency had not yet exercised that authority (\href{https://supreme.justia.com/cases/federal/us/549/497/}{Justia}). Fifteen months after the petition was filed, EPA requested public comment (66 Fed. Reg. 7486 (2001)) and received more than 50,000 comments over five months. Meanwhile, at the White House's request, the National Research Council (NRC) issued a 2001 report concluding that greenhouse gases were accumulating in Earth's atmosphere due to human activities and causing surface air and subsurface ocean temperatures to rise (\href{https://supreme.justia.com/cases/federal/us/549/497/}{Justia}).

On September 8, 2003, EPA entered an order denying the rulemaking petition (68 Fed. Reg. 52922) on two independent grounds: (1) reversing its former General Counsels, EPA concluded that the CAA does not authorize mandatory regulations to address global climate change and that greenhouse gases are therefore not ``air pollutants'' under the Act's regulatory provisions (invoking \textit{FDA v. Brown \& Williamson Tobacco Corp.}, 529 U.S. 120 (2000), post-1990 congressional inaction, and potential overlap with Department of Transportation fuel-economy standards); and (2) even if EPA possessed statutory authority, it would decline to exercise it because scientific uncertainty remained, vehicle-only standards would be an ``inefficient, piecemeal approach'' conflicting with the President's voluntary programs, and unilateral U.S. regulation might impair the President's ability to negotiate emissions reductions with developing nations such as China and India (\href{https://www.law.cornell.edu/supct/html/05-1120.ZO.html}{Cornell LII, Majority Opinion}).

A coalition of twelve States (California, Connecticut, Illinois, Maine, Massachusetts, New Jersey, New Mexico, New York, Oregon, Rhode Island, Vermont, and Washington), four local and territorial governments (District of Columbia, American Samoa, New York City, and Baltimore), and thirteen private environmental organizations petitioned for review (\href{https://supreme.justia.com/cases/federal/us/549/497/}{Justia}). Ten intervening States (Alaska, Idaho, Kansas, Michigan, Nebraska, North Dakota, Ohio, South Dakota, Texas, and Utah) and six industry trade associations intervened in support of EPA (\href{https://supreme.justia.com/cases/federal/us/549/497/}{Justia}). To establish Article III standing, petitioners submitted uncontested expert affidavits---including declarations from climate scientist Dr. Michael MacCracken, Dr. Paul H. Kirshen, and Massachusetts coastal official Karst R. Hoogeboom---showing that 20th-century sea-level rise of 10 to 20 centimeters had already begun to erode Massachusetts's substantial state-owned coastal property (including approximately 53 coastal state parks, beaches, reservations, and related infrastructure), that projected sea-level rise could permanently inundate additional state land with remediation costs running into the hundreds of millions of dollars, and that U.S. transportation emissions exceeded 1.7 billion metric tons of $\mathrm{CO_2}$ in 1999---more than 6\% of global $\mathrm{CO_2}$ emissions (\href{https://supreme.justia.com/cases/federal/us/549/497/}{Justia}).

\section{Procedural History}

Petitioners sought direct review of EPA's denial order in the United States Court of Appeals for the District of Columbia Circuit pursuant to 42 U.S.C. \S\ 7607(b)(1) (\href{https://supreme.justia.com/cases/federal/us/549/497/}{Justia}). On September 13, 2005, a sharply fractured three-judge panel denied the petition for review in \textit{Massachusetts v. EPA}, 415 F.3d 50 (D.C. Cir. 2005), with each judge writing separately:
\begin{itemize}
    \item \textbf{Judge Randolph} announced the court's judgment, declining to resolve standing definitively on the ground that standing and merits overlapped, assuming EPA possessed statutory authority, and holding that EPA reasonably exercised its ``judgment'' under \S\ 7521(a)(1) based on scientific uncertainty and foreign-policy considerations (\href{https://supreme.justia.com/cases/federal/us/549/497/}{Justia}).
    \item \textbf{Judge Sentelle} dissented in part and concurred in the judgment, concluding that petitioners lacked Article III standing because global warming is a generalized grievance ``harmful to humanity at large'' rather than a particularized injury, but accepting the contrary view as law of the case to join Judge Randolph's merits disposition (\href{https://supreme.justia.com/cases/federal/us/549/497/}{Justia}).
    \item \textbf{Judge Tatel} dissented, concluding that at least Massachusetts satisfied injury in fact, causation, and redressability based on its uncontested coastal-loss affidavits, that the CAA unambiguously covers greenhouse gases, and that EPA's policy rationales were statutorily impermissible (\href{https://supreme.justia.com/cases/federal/us/549/497/}{Justia}).
\end{itemize}

The D.C. Circuit denied rehearing en banc, 433 F.3d 66 (D.C. Cir. 2005). The Supreme Court granted certiorari on June 26, 2006 (548 U.S. 903 (2006)), heard oral argument on November 29, 2006, and issued its decision on April 2, 2007 (\href{https://supreme.justia.com/cases/federal/us/549/497/}{Justia}).

\section{Judicial Votes}

The Supreme Court decided the case by a \textbf{5--4} vote, reversing the D.C. Circuit and remanding for further proceedings (\href{https://www.law.cornell.edu/supct/html/05-1120.ZS.html}{Cornell LII, Syllabus}):
\begin{itemize}
    \item \textbf{Majority Opinion:} Justice Stevens delivered the opinion of the Court, joined by Justices Kennedy, Souter, Ginsburg, and Breyer.
    \item \textbf{Concurring Opinions:} None.
    \item \textbf{Dissenting Opinion (Standing):} Chief Justice Roberts filed a dissenting opinion, joined by Justices Scalia, Thomas, and Alito.
    \item \textbf{Dissenting Opinion (Merits):} Justice Scalia filed a dissenting opinion, joined by Chief Justice Roberts and Justices Thomas and Alito.
\end{itemize}

\section{Holding}

\textbf{Judgment:} Reversed and remanded (415 F.3d 50 reversed and remanded) (\href{https://www.law.cornell.edu/supct/html/05-1120.ZS.html}{Cornell LII, Syllabus}).

\textbf{Rule Announced:}
\begin{enumerate}
    \item \textbf{Article III Standing and ``Special Solicitude'':} At least one petitioner---here, the Commonwealth of Massachusetts---has Article III standing to challenge EPA's denial of a rulemaking petition under 42 U.S.C. \S\ 7607(b)(1). Because a State possesses quasi-sovereign interests in its territory and surrendered sovereign self-help prerogatives upon entering the Union, a State asserting a congressionally conferred procedural right to protect its concrete proprietary and quasi-sovereign interests is entitled to ``special solicitude'' in the standing analysis (\href{https://supreme.justia.com/cases/federal/us/549/497/}{Justia}). Under that framework: (a) a concrete, actual, and imminent injury (loss of state-owned coastal land from rising seas) is not rendered nonjusticiable merely because climate risks are widely shared; (b) an agency's refusal to regulate a meaningful contributor to global emissions (over 6\% of global $\mathrm{CO_2}$) satisfies traceability; and (c) redressability is satisfied where requested relief would incrementally slow or reduce the risk of harm, even if it will not single-handedly reverse global warming (\href{https://supreme.justia.com/cases/federal/us/549/497/}{Justia}).
    \item \textbf{Clean Air Act Authority:} Greenhouse gases unambiguously fall within the Clean Air Act's sweeping definition of ``air pollutant'' in 42 U.S.C. \S\ 7602(g), giving EPA statutory authority to regulate greenhouse gas emissions from new motor vehicles under \S\ 7521(a)(1) (\href{https://www.law.cornell.edu/supct/html/05-1120.ZS.html}{Cornell LII, Syllabus}).
    \item \textbf{Limits on Agency Discretion:} An agency's denial of a rulemaking petition is judicially reviewable (unlike typical nonenforcement decisions under \textit{Heckler v. Chaney}, 470 U.S. 821 (1985)). On remand, EPA may decline to regulate under \S\ 7521(a)(1) only by determining that greenhouse gases do not contribute to climate change or by offering a reasonable, statutorily grounded explanation for why it cannot or will not make that scientific judgment---not by invoking extrastatutory policy or foreign-affairs preferences (\href{https://www.law.cornell.edu/supct/html/05-1120.ZS.html}{Cornell LII, Syllabus}).
\end{enumerate}

\textbf{Holding vs.\ Dicta:} The Court's binding holdings are that Massachusetts established Article III standing (making it unnecessary to decide whether the municipal or private organizational petitioners independently had standing), that greenhouse gases are ``air pollutants'' under \S\ 7602(g) for purposes of \S\ 7521(a)(1), and that EPA's denial rested on impermissible extrastatutory grounds. By contrast, the Court expressly reserved---and did not hold---whether EPA must make an affirmative endangerment finding on remand or whether policy considerations may inform the shape of regulations once an endangerment finding is made (\href{https://supreme.justia.com/cases/federal/us/549/497/}{Justia}). Furthermore, scholars and later courts have vigorously debated whether ``special solicitude'' is an indispensable component of the standing holding or merely rhetorical dicta reinforcing Massachusetts's independent proprietary injury as a coastal landowner (\href{https://scholarship.law.umn.edu/cgi/viewcontent.cgi?article=1011&context=minnlrev}{Minnesota Law Review, 2024}; \href{https://scholarship.law.uc.edu/cgi/viewcontent.cgi?article=1550&context=uclr}{Univ.\ of Cincinnati Law Review, 2024}).

\section{Analysis of Opinions}

\subsection{Majority Opinion (Justice Stevens)}

Justice Stevens grounded the standing analysis in the constitutional baseline of \textit{Lujan v. Defenders of Wildlife}, 504 U.S. 555 (1992), while emphasizing two doctrinal amplifiers. First, Congress granted petitioners a statutory procedural right to challenge EPA's denial of a rulemaking petition in 42 U.S.C. \S\ 7607(b)(1); under \textit{Lujan} footnote 7, a litigant vested with a procedural right to protect concrete interests may assert that right ``without meeting all the normal standards for redressability and immediacy'' (\href{https://supreme.justia.com/cases/federal/us/549/497/}{Justia}). Second, because only one petitioner needs standing, Stevens focused on Massachusetts's unique status as a sovereign State. Invoking Justice Holmes's opinion in \textit{Georgia v. Tennessee Copper Co.}, 206 U.S. 230, 237 (1907), Stevens reasoned that States are ``not normal litigants for the purposes of invoking federal jurisdiction.'' Upon entering the Union, Massachusetts surrendered sovereign prerogatives to invade neighboring States, negotiate emissions treaties with China or India, or exercise preempted police powers; those powers now reside in the Federal Government, which Congress tasked EPA with exercising. Combining this quasi-sovereign territorial stake with the statutory procedural right, the Court accorded Massachusetts ``special solicitude'' in the standing inquiry (\href{https://supreme.justia.com/cases/federal/us/549/497/}{Justia}).

Applying the three Article III prongs, Stevens found: (1) \textbf{Injury in fact:} Relying on \textit{FEC v. Akins}, 524 U.S. 11, 24 (1998), widely shared climate risks do not negate a concrete, particularized harm where uncontested affidavits show rising seas have already begun swallowing Massachusetts's state-owned coastal property and threaten hundreds of millions of dollars in future loss; (2) \textbf{Causation:} EPA did not dispute that anthropogenic greenhouse gases drive global warming, and U.S. motor vehicles account for over 6\% of global $\mathrm{CO_2}$ emissions---hardly a trivial contribution; and (3) \textbf{Redressability:} Agencies, like legislatures, ``whittle away'' at massive problems incrementally (\textit{Williamson v. Lee Optical}, 348 U.S. 483, 489 (1955)), and slowing the pace of global emissions reduces Massachusetts's risk of catastrophic harm regardless of rising emissions in China or India (\href{https://supreme.justia.com/cases/federal/us/549/497/}{Justia}).

On the merits, Stevens first distinguished \textit{Heckler v. Chaney}, 470 U.S. 821 (1985), explaining that denials of rulemaking petitions are formal, legal determinations expressly made reviewable under \S\ 7607(b)(1), though review remains ``extremely limited'' and ``highly deferential.'' Turning to statutory authority, Stevens held that \S\ 7602(g)'s repeated use of ``any'' unambiguously sweeps in greenhouse gases. He distinguished \textit{FDA v. Brown \& Williamson Tobacco Corp.}, 529 U.S. 120 (2000), because CAA regulation would not compel an outright ban on motor vehicles and because EPA had never maintained a history of disclaiming authority over $\mathrm{CO_2}$. Finally, Stevens rejected EPA's alternative policy rationales (voluntary programs, piecemeal regulation, and foreign-policy leverage) as ``divorced from the statutory text'' of \S\ 7521(a)(1), holding that EPA must ground any refusal to act in the statutory endangerment inquiry itself (\href{https://www.law.cornell.edu/supct/html/05-1120.ZO.html}{Cornell LII, Majority Opinion}).

\subsection{Dissenting Opinion on Standing (Chief Justice Roberts)}

Chief Justice Roberts argued that petitioners' suit was a nonjusticiable policy dispute belonging in the elected branches, not federal court (\href{https://www.law.cornell.edu/supct/html/05-1120.ZD.html}{Cornell LII, Roberts Dissent}). He attacked the majority's ``special solicitude'' doctrine on four fronts: (1) nothing in 42 U.S.C. \S\ 7607(b)(1) grants States heightened standing compared to private litigants; (2) \textit{Georgia v. Tennessee Copper Co.} concerned equitable remedies in an original-jurisdiction \textit{parens patriae} suit, not a relaxation of Article III standing; (3) \textit{parens patriae} status is an \textit{additional} prudential hurdle rather than a substitute for Article III injury, and Massachusetts's proprietary land ownership is a ``nonsovereign interest'' under \textit{Alfred L. Snapp \& Son, Inc. v. Puerto Rico ex rel. Barez}, 458 U.S. 592, 601 (1982); and (4) \textit{Massachusetts v. Mellon}, 262 U.S. 447, 485--486 (1923), generally bars a State from suing the Federal Government as \textit{parens patriae} (\href{https://supreme.justia.com/cases/federal/us/549/497/}{Justia}).

Applying traditional \textit{Lujan} standards, Roberts found every prong lacking. Global warming is ``harmful to humanity at large'' and thus incompatible with particularized injury; petitioners' coastal-loss projections rested on a 100-year horizon and elevation models whose margin of error (30 to 70 cm) matched or exceeded projected sea-level rise, failing ``imminence.'' On causation and redressability, U.S. new motor vehicles represent only a fraction of 4\% of global greenhouse gas emissions, and any modest reduction would be overwhelmed by the 80\% of global emissions originating abroad, particularly in China and India. Roberts condemned the majority's incremental-remedy theory as a ``sleight-of-hand'' disconnecting emissions reductions from actual coastal land preservation, warning that the decision revived the universally criticized \textit{United States v. SCRAP}, 412 U.S. 669 (1973), as ``SCRAP for a new generation'' (\href{https://supreme.justia.com/cases/federal/us/549/497/}{Justia}).

\subsection{Dissenting Opinion on the Merits (Justice Scalia)}

Joining Chief Justice Roberts's standing dissent in full, Justice Scalia dissented separately on the merits (\href{https://www.law.cornell.edu/supct/html/05-1120.ZD1.html}{Cornell LII, Scalia Dissent}). First, Scalia argued that nothing in \S\ 7521(a)(1) compels the EPA Administrator to form an endangerment ``judgment'' whenever a private group files a rulemaking petition; while the statute specifies what the judgment must be about once formed, it is silent on permissible reasons for \textit{deferring} judgment, entitling EPA's policy rationales to deference under \textit{Chevron U.S.A. Inc. v. Natural Resources Defense Council, Inc.}, 467 U.S. 837 (1984). Moreover, even on the majority's own standard, EPA had extensively quoted the 2001 NRC Report's findings of substantial scientific uncertainty (\href{https://supreme.justia.com/cases/federal/us/549/497/}{Justia}).

Second, Scalia argued that EPA reasonably construed ``air pollutant'' in \S\ 7602(g) to exclude greenhouse gases. The statutory definition requires a substance to be an ``air pollution agent'' \textit{as well as} a physical or chemical substance emitted into the ambient air; reading the clause following ``including'' to swallow the preceding limitation meant that ``everything airborne, from Frisbees to flatulence, qualifies as an `air pollutant''' (\href{https://www.law.cornell.edu/supct/html/05-1120.ZD1.html}{Cornell LII, Scalia Dissent}). Because ``air pollution'' is undefined in the CAA, EPA reasonably interpreted it to mean ground-level impurities in ambient air near Earth's surface rather than globally mixed upper-atmospheric heat-trapping gases. Scalia concluded that the Court's alarm over global warming had led it to substitute its preferred policy outcome for the reasoned judgment of the responsible agency (\href{https://www.law.cornell.edu/supct/html/05-1120.ZD1.html}{Cornell LII, Scalia Dissent}).

\section{Examples: Future Applications}

\begin{enumerate}
    \item \textbf{Same-Side (Standing Likely Granted): State Challenge to Federal Fisheries Inaction Threatening State Coastal Waters.} A coastal State that owns submerged tidal lands and manages state fisheries petitions the National Marine Fisheries Service under a federal statute granting judicial review of rulemaking denials, seeking limits on offshore industrial discharges that warm coastal waters and degrade state-owned shellfish beds. Even though ocean warming has multiple global causes and federal rules would only incrementally reduce damage, the State's proprietary and quasi-sovereign territorial injury paired with a statutory procedural right closely tracks Massachusetts's posture and should establish Article III standing under \textit{Massachusetts v. EPA} (\href{https://supreme.justia.com/cases/federal/us/549/497/}{Justia}).

    \item \textbf{Same-Side (Standing Likely Granted): State Challenge to Federal Preemption or Rollback of Interstate Air-Quality Standards.} A downwind State petitions EPA to revise interstate ozone transport rules under the Clean Air Act, submitting uncontested air-monitoring data showing that upwind emissions cause its state-owned parks and residents to experience heightened ozone-attainment compliance costs and physical resource degradation. Because the State asserts both a statutory review right under \S\ 7607(b)(1) and a concrete territorial/proprietary stake surrendered to federal regulation upon joining the Union, it satisfies injury, traceability, and incremental redressability (\href{https://supreme.justia.com/cases/federal/us/549/497/}{Justia}).

    \item \textbf{Opposite-Side (Standing Likely Denied): Private Environmental Group Suing Without a State Co-Plaintiff Over Global Climate Inaction.} A private nonprofit environmental organization---without any State co-petitioner---sues a federal agency in district court alleging that its failure to adopt a broad national climate-mitigation policy harms the aesthetic and recreational interests of its members nationwide through projected sea-level rise over the next century. Without a sovereign State's ``special solicitude'' or proprietary territorial loss, and with attenuated causation and third-party international emissions complicating redressability, federal courts are far more likely to find the private plaintiffs' injury insufficiently concrete, imminent, or redressable under \textit{Lujan} and Chief Justice Roberts's warning against ``SCRAP for a new generation'' (\href{https://digitalcommons.wcl.american.edu/cgi/viewcontent.cgi?article=1927&context=aulr}{American Univ.\ Law Review, 2014}).

    \item \textbf{Opposite-Side (Standing Likely Denied): State Challenging Federal Executive Nonenforcement Priorities Based on Indirect Fiscal Costs.} Two States sue the Department of Homeland Security challenging federal guidelines that prioritize arresting certain categories of noncitizens over others, claiming the policy increases state spending on incarceration, public education, and emergency healthcare. Unlike \textit{Massachusetts v. EPA}---which involved a statutory rulemaking petition and direct physical loss of state territory---this suit targets core Article II prosecutorial discretion where federal courts have no historical tradition of ordering arrests and where indirect downstream fiscal costs do not confer cognizable Article III standing, even for States (\href{https://scholarship.law.uc.edu/cgi/viewcontent.cgi?article=1550&context=uclr}{Univ.\ of Cincinnati Law Review, 2024}; cf.\ \href{https://supreme.justia.com/cases/federal/us/602/23-235/}{\textit{FDA v. Alliance for Hippocratic Medicine}, 602 U.S. 367 (2024)}).

    \item \textbf{Fence-Sitter (Close Call): State Challenging Federal Agency Deregulation Based on Indirect Loss of State Tax Revenue or Regulatory Fees.} A State challenges a federal agency's repeal of an energy-efficiency regulation, arguing not that its physical territory will be harmed, but that lower federal standards will shift consumer purchasing patterns and reduce revenues collected by a state-run public utility authority or state-chartered instrumentality. On one hand, a direct financial injury to a state-created public instrumentality can suffice for proprietary Article III standing without even invoking ``special solicitude'' (as in \textit{Biden v. Nebraska}, 600 U.S. 477 (2023), discussed in \href{https://harvardlawreview.org/forum/vol-137/what-happened-to-traceability/}{Harvard Law Review Forum, 2024}). On the other hand, if the State alleges only general tax-base erosion or indirect macroeconomic ripple effects from third-party consumer choices, courts routinely reject such ``incidental'' revenue losses as too attenuated and non-cognizable under Article III (\href{https://scholarship.law.ufl.edu/cgi/viewcontent.cgi?article=1677&context=flr}{Florida Law Review, 2009}).
\end{enumerate}

\section{Critique}

\textbf{Originalist and Textualist Critiques:} Conservative and originalist scholars largely echo Chief Justice Roberts's and Justice Scalia's dissents. On standing, critics argue that ``special solicitude'' lacks grounding in Article III's text, Founding-era practice, or prior precedent, improperly conflating prudential \textit{parens patriae} doctrine with the ``irreducible constitutional minimum'' of \textit{Lujan} (\href{https://scholarship.law.ufl.edu/cgi/viewcontent.cgi?article=1677&context=flr}{Calvin Massey, 61 Fla.\ L.\ Rev.\ 249 (2009)}; \href{https://scholarlycommons.law.case.edu/cgi/viewcontent.cgi?article=1185&context=faculty_publications}{Jonathan H. Adler, Case Western Reserve Faculty Pubs.}). By allowing States to bypass \textit{Mellon}'s bar and relax imminence and redressability, \textit{Massachusetts v. EPA} invited partisan ``state attorney general litigation''---first by blue States against Republican administrations and later by red States against Democratic administrations (as seen in immigration and student-loan suits), turning federal courts into referees of national policy disputes (\href{https://scholarship.law.uc.edu/cgi/viewcontent.cgi?article=1368&context=fac_pubs}{Bradford C. Mank, Univ.\ of Cincinnati Scholarship}; \href{https://scholarship.law.uc.edu/cgi/viewcontent.cgi?article=1550&context=uclr}{Univ.\ of Cincinnati Law Review, 2024}). On the merits, textualists and administrative-law critics contend that the majority short-circuited \textit{Chevron} deference and stretched a 1970 localized air-pollution statute to cover global climate change---a tension later reflected when the Court curtailed EPA's greenhouse-gas permitting authority in \href{https://supreme.justia.com/cases/federal/us/573/302/}{\textit{Utility Air Regulatory Group v. EPA}, 573 U.S. 302 (2014)}, applied the major questions doctrine in \textit{West Virginia v. EPA} (2022), overruled \textit{Chevron} in \href{https://supreme.justia.com/cases/federal/us/603/22-451/}{\textit{Loper Bright Enterprises v. Raimondo}, 603 U.S. 369 (2024)}, and when EPA ultimately moved to rescind the 2009 Endangerment Finding itself (\href{https://www.federalregister.gov/documents/2026/02/18/2026-03157/rescission-of-the-greenhouse-gas-endangerment-finding-and-motor-vehicle-greenhouse-gas-emission}{91 Fed.\ Reg.\ 7686 (Feb.\ 18, 2026)}).

\textbf{Progressive and Doctrinal Critiques:} From a progressive and doctrinal-coherence perspective, scholars praise \textit{Massachusetts v. EPA} for preventing the Executive Branch from nullifying statutory mandates through pretextual policy excuses and for recognizing that collective, global harms cannot be immunized from judicial review simply because everyone is injured (\href{https://www.scotusblog.com/2007/04/discussion-board-thoughts-on-mass-v-epa/}{SCOTUSblog, 2007}). Yet even sympathetic scholars fault the opinion's doctrinal ambiguity. Commentators observe that the Court never defined what ``special solicitude'' actually does---whether it lowers the bar for injury, causation, redressability, or all three---leaving lower courts without a workable rule and prompting the modern Supreme Court to quietly sideline ``special solicitude'' without overruling it (\href{https://scholarship.law.umn.edu/cgi/viewcontent.cgi?article=1011&context=minnlrev}{Minnesota Law Review, 2024}; \href{https://harvardlawreview.org/forum/vol-137/what-happened-to-traceability/}{Harvard Law Review Forum, 2024}). Others highlight what they call ``the standing doctrine's dirty little secret'': if Article III's imminence and redressability limits are truly constitutional floors, the Court never satisfactorily explained why Congress's creation of a ``procedural right'' or a plaintiff's status as a State can constitutionally relax them (\href{https://scholarlycommons.law.northwestern.edu/cgi/viewcontent.cgi?article=1084&context=nulr}{Northwestern Univ.\ Law Review, 2012}).

\textbf{Underlying Competing Values:} At bottom, \textit{Massachusetts v. EPA} exposes a deep structural tension in constitutional law between two competing visions of the separation of powers and federalism. The majority prioritizes \textbf{statutory fidelity, executive accountability, and federalist compensation}: when Congress enacts broad protective statutes and preempts state self-help, courts must remain open---especially to sovereign States---so that the Executive cannot effectively repeal congressional mandates by inaction. The dissents prioritize \textbf{judicial restraint, democratic accountability, and structural limits on Article III}: global, polycentric problems like climate change involve complex trade-offs across domestic economics and foreign diplomacy that belong to Congress and the President, not to unelected judges adjudicating attenuated causal chains at the behest of dissatisfied litigants.

\section{Key Quotations}

\begin{quote}
``Given that procedural right and Massachusetts' stake in protecting its quasi-sovereign interests, the Commonwealth is entitled to special solicitude in our standing analysis.''
\end{quote}
\noindent --- Justice Stevens, Opinion of the Court, 549 U.S.\ at 520 (\href{https://supreme.justia.com/cases/federal/us/549/497/}{Justia}; \href{https://www.law.cornell.edu/supct/html/05-1120.ZO.html}{Cornell LII}).

\begin{quote}
``On its face, the definition embraces all airborne compounds of whatever stripe, and underscores that intent through the repeated use of the word `any.' Carbon dioxide, methane, nitrous oxide, and hydrofluorocarbons are without a doubt `physical [and] chemical . . . substance[s] which [are] emitted into . . . the ambient air.' The statute is unambiguous.''
\end{quote}
\noindent --- Justice Stevens, Opinion of the Court, 549 U.S.\ at 528--529 (\href{https://supreme.justia.com/cases/federal/us/549/497/}{Justia}; \href{https://www.law.cornell.edu/supct/html/05-1120.ZO.html}{Cornell LII}).

\begin{quote}
``Relaxing Article III standing requirements because asserted injuries are pressed by a State, however, has no basis in our jurisprudence, and support for any such `special solicitude' is conspicuously absent from the Court's opinion.''
\end{quote}
\noindent --- Chief Justice Roberts, dissenting, 549 U.S.\ at 536 (\href{https://supreme.justia.com/cases/federal/us/549/497/}{Justia}; \href{https://www.law.cornell.edu/supct/html/05-1120.ZD.html}{Cornell LII}).

\begin{quote}
``SCRAP made standing seem a lawyer's game, rather than a fundamental limitation ensuring that courts function as courts and not intrude on the politically accountable branches. Today's decision is SCRAP for a new generation.''
\end{quote}
\noindent --- Chief Justice Roberts, dissenting, 549 U.S.\ at 547--548 (\href{https://supreme.justia.com/cases/federal/us/549/497/}{Justia}; \href{https://www.law.cornell.edu/supct/html/05-1120.ZD.html}{Cornell LII}).

\begin{quote}
``It follows that everything airborne, from Frisbees to flatulence, qualifies as an `air pollutant.' This reading of the statute defies common sense.''
\end{quote}
\noindent --- Justice Scalia, dissenting, 549 U.S.\ at 558 n.2 (\href{https://supreme.justia.com/cases/federal/us/549/497/}{Justia}; \href{https://www.law.cornell.edu/supct/html/05-1120.ZD1.html}{Cornell LII}).

\end{document}
